% !TeX program = pdflatex
\documentclass[aspectratio=169,xcolor={dvipsnames,table}]{beamer}

\input{../../_en_preambulo_beamer}
% Comandos y macros estadísticas compartidas para las presentaciones Beamer en inglés (Metropolis)

% Conjuntos numéricos básicos
\newcommand{\R}{\mathbb{R}}
\newcommand{\N}{\mathbb{N}}
\newcommand{\Q}{\mathbb{Q}}
\newcommand{\Z}{\mathbb{Z}}
\newcommand{\set}[1]{\left\{ #1 \right\}}
\newcommand{\abs}[1]{\left|#1\right|}
\newcommand{\norm}[1]{\left\| #1 \right\|}

% Comandos estadísticos y probabilísticos del curso
\newcommand{\Var}{\operatorname{Var}}
\newcommand{\cov}{\operatorname{Cov}}
\newcommand{\corr}{\rho}
\newcommand{\comb}[2]{\begin{pmatrix} #1 \\ #2 \end{pmatrix}}
\newcommand{\s}{\sigma}
\newcommand{\particion}{\mathfrak{P}}
\newcommand{\card}[1]{\#\left( #1 \right)}
\newcommand{\seq}[3]{\set{#1}_{#2}^{#3}}
\newcommand{\E}{\mathbb{E}}
\newcommand{\Prob}{\mathbb{P}}

% Entornos pedagógicos y cajas en inglés académico natural (soportando tanto nombres en inglés como en español)
\newenvironment{discussion}{%
  \begin{alertblock}{Discussion Question}%
}{%
  \end{alertblock}%
}
\newenvironment{discusion}{%
  \begin{alertblock}{Discussion Question}%
}{%
  \end{alertblock}%
}

\newenvironment{reflection}{%
  \begin{alertblock}{Classroom Reflection}%
}{%
  \end{alertblock}%
}
\newenvironment{reflexion}{%
  \begin{alertblock}{Classroom Reflection}%
}{%
  \end{alertblock}%
}

\newenvironment{paradox}{%
  \begin{alertblock}{Exploratory Paradox}%
}{%
  \end{alertblock}%
}
\newenvironment{paradoja}{%
  \begin{alertblock}{Exploratory Paradox}%
}{%
  \end{alertblock}%
}

\newenvironment{motivation}{%
  \begin{exampleblock}{Motivation: Data Science in Practice}%
}{%
  \end{exampleblock}%
}
\newenvironment{motivacion}{%
  \begin{exampleblock}{Motivation: Data Science in Practice}%
}{%
  \end{exampleblock}%
}

\newenvironment{quickchallenge}{%
  \begin{block}{Quick Team Challenge}%
}{%
  \end{block}%
}
\newenvironment{retoclase}{%
  \begin{block}{Quick Team Challenge}%
}{%
  \end{block}%
}


\title{Geometric and Negative Binomial Distributions}
\subtitle{Section 03.05 --- Waiting Times, Memoryless Property, and Overdispersion}
\author[J. Castillo Colmenares]{Juliho Castillo Colmenares}
\institute[Tec de Monterrey]{Tecnológico de Monterrey}
\date{\vspace{-1.2cm}}

\begin{document}

% SLIDE 01: Title Page
\begin{frame}[plain]
	\titlepage
\end{frame}

% SLIDE 02: Roadmap
\begin{frame}{Roadmap --- Chapter 03: Discrete Random Variables}
	\begin{block}{Curricular Structure of Chapter 03}
		\small
		\begin{enumerate}
			\itemsep=0.04cm
			\item \textcolor{gray}{Section 03.01: Discrete Probability Mass Functions (PMF and Support)}
			\item \textcolor{gray}{Section 03.02: Cumulative Distribution Functions (Discrete CDF)}
			\item \textcolor{gray}{Section 03.03: Mathematical Expectation, Variance, and Moments}
			\item \textcolor{gray}{Section 03.04: Bernoulli and Binomial Distributions}
			\item \textbf{\textcolor{TecRojo}{Section 03.05: Geometric and Negative Binomial Distributions}} $\leftarrow$ \emph{Current focus}
			\item \textcolor{gray}{Section 03.06: Hypergeometric Distribution}
		\end{enumerate}
	\end{block}
	\vspace{-0.15cm}
	\begin{alertblock}{Learning Objective}
		\footnotesize
		Analyze the modeling of waiting times until the $r$-th success in independent Bernoulli trials, mastering the memoryless property for $r=1$ and the overdispersion phenomenon ($\sigma^2 > \mu$) for $r \ge 1$.
	\end{alertblock}
\end{frame}

% SLIDE 03: Motivation and Intuition
\begin{frame}{Motivation --- From Counting Successes to Measuring Waiting Times}
	\begin{columns}[T]
		\begin{column}{0.48\textwidth}
			\begin{block}{The Binomial Perspective (Fixed Trials)}
				\small
				\begin{itemize}
					\itemsep=0.06cm
					\item We fix beforehand the sample size or total number of trials ($n$).
					\item The random variable is the number of observed successes ($X \in \set{0, 1, \dots, n}$).
					\item Typical query: \emph{How many defective packages will appear in a batch of 100?}
				\end{itemize}
			\end{block}
		\end{column}
		\begin{column}{0.48\textwidth}
			\begin{alertblock}{The Waiting Time Perspective (Fixed Successes)}
				\small
				\begin{itemize}
					\itemsep=0.06cm
					\item We fix beforehand the number of target successes required ($r$).
					\item The random variable is the total number of trials needed ($X \in \set{r, r+1, \dots}$).
					\item Typical query: \emph{How many attempts are needed until detecting the first failure ($r=1$) or the fourth ($r=4$)?}
				\end{itemize}
			\end{alertblock}
		\end{column}
	\end{columns}
\end{frame}

% SLIDE 04: The Geometric Distribution
\begin{frame}{The Geometric Distribution ($X \sim \text{Geom}(p)$)}
	\begin{block}{Formal Definition --- Time Until the First Success}
		\small
		Let $X$ be the total number of i.i.d. Bernoulli trials required to obtain the **first success** with probability $p \in (0, 1)$ and failure probability $q = 1-p$. Its probability mass function (PMF) and support are:
		\begin{align*}
			f_X(k) = P(X = k) = q^{k-1}p = (1-p)^{k-1}p, \quad k \in \mathcal{S}_X = \set{1, 2, 3, \dots}.
		\end{align*}
	\end{block}
	\vspace{-0.15cm}
	\begin{columns}[T]
		\begin{column}{0.48\textwidth}
			\begin{exampleblock}{Exact Normalization}
				\footnotesize
				By the sum of infinite geometric series with ratio $q \in (0, 1)$:
				$\sum_{k=1}^{\infty} q^{k-1}p = p \sum_{j=0}^{\infty} q^j = p \left(\frac{1}{1-q}\right) = \frac{p}{p} = 1$.
			\end{exampleblock}
		\end{column}
		\begin{column}{0.48\textwidth}
			\begin{alertblock}{Cumulative Function (CDF)}
				\footnotesize
				The upper tail probability is $P(X > k) = q^k$. Thus, for any integer $k \ge 1$:
				$F_X(k) = P(X \le k) = 1 - q^k$.
			\end{alertblock}
		\end{column}
	\end{columns}
\end{frame}

% SLIDE 05: Geometric Moments
\begin{frame}{Moments of the Geometric Distribution}
	\begin{block}{Mathematical Expectation ($\mu$) and Variance ($\sigma^2$)}
		\small
		Using power series differentiation or the moment generating function $M_X(t) = \frac{pe^t}{1-qe^t}$:
		\begin{align*}
			\mathbb{E}[X] = \dfrac{1}{p}, \qquad \text{Var}(X) = \dfrac{q}{p^2} = \dfrac{1-p}{p^2}, \qquad \sigma_X = \dfrac{\sqrt{1-p}}{p}.
		\end{align*}
	\end{block}
	\vspace{-0.15cm}
	\begin{columns}[T]
		\begin{column}{0.48\textwidth}
			\begin{exampleblock}{Mean Interpretation}
				\footnotesize
				If the probability of success in a single attempt is $p = 0.20$ (1 in 5), on average it takes $\mathbb{E}[X] = \frac{1}{0.20} = 5$ attempts to achieve the first success.
			\end{exampleblock}
		\end{column}
		\begin{column}{0.48\textwidth}
			\begin{alertblock}{Positive Skewness}
				\footnotesize
				The mode is always exactly $k=1$ ($f_X(1)=p > qp = f_X(2)$). The distribution is strictly decreasing with a heavy right tail.
			\end{alertblock}
		\end{column}
	\end{columns}
\end{frame}

% SLIDE 06: Memoryless Property
\begin{frame}{The Memoryless Property (Discrete Markovianity)}
	\begin{block}{Time-Invariance Theorem ($X \sim \text{Geom}(p)$)}
		\small
		For any pair of positive integers $m, n \in \set{1, 2, \dots}$, the conditional probability of requiring $n$ additional trials given survival past the first $m$ failed attempts satisfies:
		\begin{align*}
			P(X > m+n \mid X > m) = P(X > n) = q^n.
		\end{align*}
	\end{block}
	\vspace{-0.15cm}
	\begin{columns}[T]
		\begin{column}{0.48\textwidth}
			\begin{exampleblock}{Proof via Tail Probabilities}
				\scriptsize
				By conditional definition for $m, n \ge 1$:
				\begin{align*}
					P(X > m+n \mid X > m) &= \dfrac{P(X > m+n)}{P(X > m)} \\
					&= \dfrac{q^{m+n}}{q^m} = q^n.
				\end{align*}
			\end{exampleblock}
		\end{column}
		\begin{column}{0.48\textwidth}
			\begin{alertblock}{Unique Characterization}
				\footnotesize
				The geometric distribution is the **only** discrete law on $\set{1, 2, \dots}$ lacking memory. The system "does not age" nor remember past failures.
			\end{alertblock}
		\end{column}
	\end{columns}
\end{frame}

% SLIDE 07: The Negative Binomial Distribution
\begin{frame}{The Negative Binomial Distribution ($X \sim \text{BN}(r, p)$)}
	\begin{block}{Combinatorial Generalization --- Time Until the $r$-th Success}
		\small
		Let $X$ be the total number of trials required to accumulate **exactly $r$ successes** ($r \ge 1$). For the $r$-th success to happen on trial $k$, there must be $(r-1)$ successes in the first $(k-1)$ trials combined with success on trial $k$:
		\begin{align*}
			f_X(k) = P(X = k) = \comb{k-1}{r-1} p^r q^{k-r}, \quad k \in \mathcal{S}_X = \set{r, r+1, r+2, \dots}.
		\end{align*}
	\end{block}
	\vspace{-0.15cm}
	\begin{columns}[T]
		\begin{column}{0.48\textwidth}
			\begin{exampleblock}{Pascal's Relation}
				\footnotesize
				When $r=1$, the combinatorial coefficient reduces to $\comb{k-1}{0} = 1$, exactly recovering the elementary Geometric law $P(X=k)=q^{k-1}p$.
			\end{exampleblock}
		\end{column}
		\begin{column}{0.48\textwidth}
			\begin{alertblock}{Alternative Parameterization}
				\footnotesize
				In statistical software (`scipy`), it is modeled via failure counts $Y = X - r \in \set{0, 1, \dots}$ with PMF $P(Y=y) = \comb{y+r-1}{y}p^r q^y$.
			\end{alertblock}
		\end{column}
	\end{columns}
\end{frame}

% SLIDE 08: Negative Binomial Moments
\begin{frame}{Negative Binomial Moments and Stochastic Additivity}
	\begin{block}{Decomposition as a Sum of Independent Geometric RVs}
		\small
		The total time $X$ until the $r$-th success is the stochastic sum of independent inter-success waiting times $Y_i \sim \text{Geom}(p)$: $X = \sum_{i=1}^r Y_i$. By linearity of expectation and variance:
		\begin{align*}
			\mathbb{E}[X] = \sum_{i=1}^r \mathbb{E}[Y_i] = \dfrac{r}{p}, \qquad \text{Var}(X) = \sum_{i=1}^r \text{Var}(Y_i) = \dfrac{rq}{p^2} = \dfrac{r(1-p)}{p^2}.
		\end{align*}
	\end{block}
	\vspace{-0.15cm}
	\begin{columns}[T]
		\begin{column}{0.48\textwidth}
			\begin{exampleblock}{Moment Generating Function}
				\footnotesize
				By independence, the MGF is the exact product:
				$M_X(t) = \prod_{i=1}^r M_{Y_i}(t) = \left(\frac{pe^t}{1-qe^t}\right)^r$ for $t < -\ln(q)$.
			\end{exampleblock}
		\end{column}
		\begin{column}{0.48\textwidth}
			\begin{alertblock}{Standard Deviation}
				\footnotesize
				The standard deviation of the total counting process is $\sigma_X = \frac{\sqrt{r(1-p)}}{p}$, scaling as $\sqrt{r}$ with increasing success target $r$.
			\end{alertblock}
		\end{column}
	\end{columns}
\end{frame}

% SLIDE 09: Overdispersion Analysis
\begin{frame}{Overdispersion Analysis --- Comparison of Counting Models}
	\begin{block}{The Variance-to-Mean Ratio in Discrete Processes}
		\small
		When comparing the three major discrete counting distributions for observed failures ($Y \in \set{0, 1, \dots}$), the dispersion index $\mathcal{I} = \frac{\sigma^2}{\mu}$ dictates the structural fit:
	\end{block}
	\vspace{-0.1cm}
	\begin{columns}[T]
		\begin{column}{0.32\textwidth}
			\begin{exampleblock}{Binomial ($\mathcal{I} < 1$)}
				\footnotesize
				$\mu = np, \; \sigma^2 = npq$.
				
				Ratio $\frac{\sigma^2}{\mu} = q < 1$.
				
				**Underdispersion:** Bounded trial processes with controlled variability.
			\end{exampleblock}
		\end{column}
		\begin{column}{0.32\textwidth}
			\begin{block}{Poisson ($\mathcal{I} = 1$)}
				\footnotesize
				$\mu = \lambda, \; \sigma^2 = \lambda$.
				
				Ratio $\frac{\sigma^2}{\mu} = 1$.
				
				**Equidispersion:** Rare events occurring homogeneously in continuous time.
			\end{block}
		\end{column}
		\begin{column}{0.32\textwidth}
			\begin{alertblock}{Neg. Binomial ($\mathcal{I} > 1$)}
				\footnotesize
				$\mu = \frac{rq}{p}, \; \sigma^2 = \frac{rq}{p^2}$.
				
				Ratio $\frac{\sigma^2}{\mu} = \frac{1}{p} > 1$.
				
				**Overdispersion:** Clustered arrivals or unobserved individual heterogeneity.
			\end{alertblock}
		\end{column}
	\end{columns}
\end{frame}

% SLIDE 10: Asymptotic Relation
\begin{frame}{Asymptotic Relation --- From Pascal to Poisson}
	\begin{block}{Convergence of Failure Counts ($r \to \infty, p \to 1$)}
		\small
		If the required successes grow indefinitely ($r \to \infty$) while single-trial probability approaches certainty ($p \to 1$), holding the mean failure rate $(1-p)r = \lambda > 0$ constant:
		\begin{align*}
			\lim_{r \to \infty} P(Y = y) = \lim_{r \to \infty} \comb{y+r-1}{y} p^r (1-p)^y = \dfrac{\lambda^y e^{-\lambda}}{y!}, \quad \forall y \in \set{0, 1, \dots}.
		\end{align*}
	\end{block}
	\vspace{-0.15cm}
	\begin{columns}[T]
		\begin{column}{0.48\textwidth}
			\begin{exampleblock}{Combinatorial Mechanism}
				\footnotesize
				The combinatorial factor $\frac{r(r+1)\cdots(r+y-1)}{y! r^y} \to \frac{1}{y!}$, while the probability term $\left(1 - \frac{\lambda}{r}\right)^r \to e^{-\lambda}$.
			\end{exampleblock}
		\end{column}
		\begin{column}{0.48\textwidth}
			\begin{alertblock}{Theoretical Synthesis}
				\footnotesize
				The Negative Binomial absorbs extra variance; as dispersion diminishes ($\frac{1}{p} \to 1$), it collapses exactly into the Poisson distribution.
			\end{alertblock}
		\end{column}
	\end{columns}
\end{frame}









% SLIDE 17: Python Lab --- PMF/CDF Validation
\begin{frame}[fragile]{Python Lab --- Vectorized PMF/CDF Validation (`scipy.stats`)}
	\vspace{-0.15cm}
	\begin{block}{Implementation in \texttt{03.05\_geometric\_negative\_binomial.py}}
		\lstinputlisting[language=Python, firstline=16, lastline=37, basicstyle=\fontsize{5.8pt}{6.8pt}\ttfamily]{../../code/03_variables_aleatorias_discretas/03.05_geometric_negative_binomial.py}
	\end{block}
\end{frame}

% SLIDE 18: Python Lab --- Memoryless Property
\begin{frame}[fragile]{Python Lab --- Memoryless Property via Monte Carlo ($N=250,000$)}
	\begin{block}{Empirical Simulation of Conditional Invariance}
		\lstinputlisting[language=Python, firstline=42, lastline=58, basicstyle=\fontsize{6.3pt}{7.5pt}\ttfamily]{../../code/03_variables_aleatorias_discretas/03.05_geometric_negative_binomial.py}
	\end{block}
\end{frame}

% SLIDE 19: Python Lab --- Overdispersion
\begin{frame}[fragile]{Python Lab --- Overdispersion and Stochastic Sum (`numpy`)}
	\begin{block}{Verification of Additivity and Variance-to-Mean Ratio ($>1.0$)}
		\lstinputlisting[language=Python, firstline=63, lastline=77, basicstyle=\fontsize{6.3pt}{7.5pt}\ttfamily]{../../code/03_variables_aleatorias_discretas/03.05_geometric_negative_binomial.py}
	\end{block}
\end{frame}

% SLIDE 20: Closing and Stochastic Synthesis
\begin{frame}{Closing --- Stochastic Synthesis and Didactic Bridge}
	\begin{columns}[T]
		\begin{column}{0.48\textwidth}
			\begin{block}{Conceptual Summary}
				\small
				\begin{itemize}
					\itemsep=0.08cm
					\item **Geometric ($p$):** Time to 1st success, $\mu = \frac{1}{p}$, $\sigma^2 = \frac{q}{p^2}$, unique memoryless discrete law.
					\item **Negative Binomial ($r, p$):** Sum of $r$ independent geometrics, $\mu = \frac{r}{p}$, natural **overdispersion** absorber ($\sigma^2 > \mu$).
					\item **Limit:** Converges to Poisson as $r \to \infty$.
				\end{itemize}
			\end{block}
		\end{column}
		\begin{column}{0.48\textwidth}
			\begin{alertblock}{Bridge to Section 03.06}
				\small
				So far, all Bernoulli trials assumed mutual independence (**sampling with replacement** from finite populations).
				
				What happens when we draw samples **without replacement** from a finite population partitioned into two distinct sub-populations?
				
				$\to$ This leads naturally to the **Hypergeometric Distribution**.
			\end{alertblock}
		\end{column}
	\end{columns}
\end{frame}

\end{document}
